\chapter{Systemarchitektur}
\label{kap:architektur}
\verfasser{\autorDG}

\section{Gesamtübersicht}
\label{sec:gesamt}

Abbildung~\ref{fig:architektur} zeigt den Aufbau des Systems. Es besteht aus zwei Teilen: einem React-Client, der im Browser läuft, und einem Node.js-Prozess als Server.

Der Client ist eine \gls{spa}, die mit Vite und TypeScript gebaut wird. Er gliedert sich in drei Schichten. Die Seiten-Komponenten bilden die Oberfläche. Der Zustand-Store hält alle Daten, die die Oberfläche anzeigt. Der Hook \texttt{useWebSocket} verwaltet die Verbindung zum Server, übersetzt Nutzeraktionen in Nachrichten und schreibt eingehende Nachrichten in den Store.

Im Server-Prozess laufen Express und der WebSocket-Server der Bibliothek \texttt{ws} auf demselben \gls{http}-Server (Port~4000). Express beantwortet die \gls{rest}-Anfragen zur Registrierung, Anmeldung, Erneuerung der Sitzung und Abmeldung und setzt dabei die Session-Cookies. Der WebSocket-Server nimmt Verbindungen auf dem Pfad \texttt{/ws} an, sofern der Handshake ein gültiges Access-Token enthält, und leitet eingehende Nachrichten an Handler weiter. Die Handler lesen und verändern den Zustand im Speicher (Lobbys, verbundene Spieler, Einladungen) und greifen über das Modul \texttt{database.js} auf die SQLite-Datenbank zu. Anschließend verteilen sie das Ergebnis an die betroffenen Clients. Datenbankzugriffe und das Hashen von Passwörtern laufen in Worker Threads (Abschnitt~\ref{sec:worker}).


Der Aufbau entspricht der klassischen Vierschichtenarchitektur für Webanwendungen \X[Quelle zur Vierschichtenarchitektur]: Die React-Komponenten bilden die Visualisierung, Store und Hook im Browser sowie die \gls{rest}- und WebSocket-Schnittstellen des Servers die Präsentation, die Handler mit dem Zustand im Speicher die Anwendungslogik und die SQLite-Datenbank die Datenhaltung.

\begin{figure}[ht]
\centering
\resizebox{\linewidth}{!}{%
\begin{tikzpicture}[
  font=\sffamily\small, >=Stealth,
  box/.style={draw, rounded corners=2pt, fill=white, minimum width=3.6cm, minimum height=1cm, align=center},
  sbox/.style={box, minimum width=3cm},
  grp/.style={draw, rounded corners=4pt, fill=black!4, inner sep=10pt},
  sub/.style={draw, dashed, rounded corners=4pt, fill=white, inner sep=7pt},
  lbl/.style={font=\sffamily\scriptsize, align=center, fill=white, inner sep=1.5pt},
  ilbl/.style={font=\sffamily\scriptsize, align=left, inner sep=1.5pt},
  slbl/.style={font=\sffamily\scriptsize\bfseries}
]
  % Client
  \node[box] (ui)    at (2.6,4.2) {Seiten-Komponenten\\(Auth, Hauptmenü,\\Lobby, Chat)};
  \node[box, minimum width=2.6cm] (store) at (2.1,2.0) {Zustand-Store};
  \node[box] (hook)  at (2.6,0)   {\texttt{useWebSocket}-Hook};
  \coordinate (cpad) at (5.6,0);

  % Server: Haupt-Thread
  \node[box]  (express) at (10.6,4.2) {Express\\\texttt{/login}, \texttt{/register}, \texttt{/auth/*}};
  \node[box]  (wss)     at (10.6,0)   {\texttt{ws}-Server\\Pfad \texttt{/ws}};
  \node[sbox] (handler) at (15,0)     {Router und Handler\\+ Broadcasting};
  \node[sbox] (mem)     at (15,2.2)   {In-Memory-Zustand\\(Lobbys, Spieler,\\Einladungen)};
  % Server: Worker Threads
  \node[sbox] (dbstore) at (15,-3.4)  {Player-, Friends-,\\Refresh-Token-Store};
  \node[cylinder, shape border rotate=90, draw, fill=white, aspect=0.25,
        minimum height=1.3cm, minimum width=2.2cm] (db) at (10.6,-3.4) {SQLite};
  \node[sbox] (pw) at (20.2,4.2) {bcrypt\\(Passwort-Hash)};
  \coordinate (spad) at (21.9,0);

  \begin{scope}[on background layer]
    \node[grp, fit=(ui)(store)(hook)(cpad), label={[font=\sffamily\small\bfseries]above:Browser: React-Client}] (client) {};
    % outer group first, the thread groups are drawn on top of it
    \node[grp, fit=(express)(wss)(handler)(mem)(dbstore)(db)(pw)(spad), inner sep=24pt, label={[font=\sffamily\small\bfseries]above:Node.js-Prozess (HTTP-Server, Port 4000)}] (server) {};
    \node[sub, fit=(express)(wss)(handler)(mem), label={[slbl]above:Haupt-Thread (Event Loop)}] (main) {};
    \node[sub, fit=(dbstore)(db), label={[slbl]below:Datenbank-Worker (1 Worker Thread)}] (dbw) {};
    \node[sub, fit=(pw), label={[slbl]above:Passwort-Worker (Pool)}] (pww) {};
  \end{scope}

  % Client intern
  \draw[->] (1.4,2.5) -- node[ilbl, right] {Re-Rendering} (1.4,3.45);
  \draw[->] (1.4,0.5) -- node[ilbl, right] {\texttt{getState()}} (1.4,1.5);
  \draw[->] (4.1,3.45) -- node[ilbl, right, pos=0.5] {Aktionen\\(\texttt{joinLobby},\\\texttt{sendMessage}, \ldots)} (4.1,0.5);

  % Netz
  \draw[->] (ui) -- node[lbl, above] {HTTP POST (Login)} node[lbl, below] {Antwort: Cookies} (express);
  \draw[<->, line width=2.2pt, black!70] (hook) -- node[lbl, above=3pt] {\textbf{WebSocket-Verbindung}} node[lbl, below=3pt] {JSON-Nachrichten mit \texttt{type}} (wss);

  % Server intern
  \draw[<->] (wss) -- (handler);
  \draw[<->] (handler) -- (mem);
  \draw[<->] (handler) -- node[lbl, right] {\texttt{postMessage}\\/ Promise} (dbstore);
  \draw[<->] (dbstore) -- (db);
  \draw[<->] (express) -- node[lbl, above] {\texttt{postMessage} / Promise} (pw);
  \draw[->] ($(express.east)+(0,-0.3)$) -| (18.2,-3.2) -- ($(dbstore.east)+(0,0.2)$);
\end{tikzpicture}}
\caption{Gesamtarchitektur des Systems}
\label{fig:architektur}
\end{figure}

\section{Aufteilung zwischen REST und WebSocket}
\label{sec:restws}

Das System verwendet zwei Kommunikationswege. Alles, was die Anmeldung betrifft, läuft über \gls{rest}-Endpunkte von Express:
\begin{itemize}
  \item \texttt{POST /register} und \texttt{POST /login} nehmen Benutzername und Passwort als \gls{json} entgegen und starten eine Sitzung.
  \item \texttt{POST /auth/refresh} erneuert eine Sitzung.
  \item \texttt{POST /auth/logout} beendet sie.
\end{itemize}
Der Server stellt dabei ein kurzlebiges Access-Token (\gls{jwt}) \X[RFC 7519] und ein langlebiges Refresh-Token aus. Beide setzt er mit dem Header \texttt{Set-Cookie} als Cookies, in Antwortkörpern sind sie nie enthalten (Abschnitt~\ref{sec:wsauth}). Alle weiteren Interaktionen laufen über die WebSocket-Verbindung, darunter Lobbys, Chat, Freunde und Einladungen.

Die Aufteilung folgt der Art der Kommunikation. Anmeldung, Erneuerung und Abmeldung sind einmalige Vorgänge mit genau einer Anfrage und einer Antwort, für die das Request-Response-Modell von \gls{http} ausreicht. Außerdem lassen sich Cookies nur in \gls{http}-Antworten setzen, nicht über eine WebSocket-Verbindung. Zudem existiert vor der Anmeldung noch keine Identität, an die eine WebSocket-Verbindung gebunden werden könnte.

Bei allen übrigen Funktionen muss eine Änderung sofort bei anderen Nutzern ankommen, etwa eine Chatnachricht bei allen Mitgliedern einer Lobby. Der Server muss diesen Nutzern also Nachrichten senden, ohne dass sie vorher angefragt haben. Das ist nur über die WebSocket-Verbindung möglich. Die Aktionen des Nutzers selbst, etwa das Senden der Chatnachricht, laufen über dieselbe Verbindung.

\section{Nachrichtenprotokoll}
\label{sec:protokoll}

Alle Nachrichten über die WebSocket-Verbindung sind \gls{json}-Objekte mit einem Feld \texttt{type}, das die Art der Nachricht angibt. Die übrigen Felder hängen vom Typ ab. Die Typen sind auf beiden Seiten als Konstanten definiert: im Server in \texttt{messageTypes.js}, im Client in \texttt{structs.tsx}. Die Konstanten für Nachrichten vom Client an den Server beginnen mit \texttt{Request}, die für Nachrichten vom Server an den Client mit \texttt{Response}. Listing~\ref{lst:protokoll-bsp} zeigt eine Anfrage zum Betreten einer Lobby und die daraufhin verschickten Nachrichten.

\begin{lstlisting}[language=JavaScript, caption={Beispiel: Betreten einer Lobby (IDs gekürzt)}, label=lst:protokoll-bsp, numbers=none]
// Client -> Server
{ "type": "join_lobby", "lobbyID": "2ccc71..", "password": "" }

// Server -> beitretender Client
{ "type": "join_lobby_successful",
  "lobby": { "id": "2ccc71..", "name": "Abendrunde", "maxPlayers": 4,
             "isPrivate": false, "players": [ ... ], "gameStart": [ ... ] },
  "chatHistory": [ ... ] }

// Server -> alle Mitglieder der Lobby
{ "type": "lobby_updated", "lobby": { ... } }

// Server -> alle verbundenen Clients
{ "type": "lobby_list", "lobbies": [ ... ] }
\end{lstlisting}

Sobald die Verbindung aufgebaut ist, sendet der Server von sich aus eine Nachricht vom Typ \texttt{welcome}. Sie enthält den angemeldeten Spieler, die Liste aller Lobbys, die offenen Freundschaftsanfragen und die Freundesliste mit Online-Status. Mit diesem Schnappschuss kann der Client seine Oberfläche vollständig aufbauen. Danach halten einzelne Ereignisnachrichten ihn aktuell. Anhang~\ref{anh:protokoll} listet alle Nachrichtentypen auf.

\section{Authentifizierung der WebSocket-Verbindung}
\label{sec:wsauth}

\paragraph{Anforderungen.} Jede WebSocket-Verbindung muss einem angemeldeten Spieler zugeordnet werden. Das Token, das diese Zuordnung belegt, darf dabei weder für Skripte auf der Seite lesbar sein noch von fremden Webseiten verwendet werden können. Da die Browser-\gls{api} beim Aufbau einer WebSocket-Verbindung keine eigenen Header zulässt (Abschnitt~\ref{sec:websocket}), scheidet ein \texttt{Authorization}-Header aus.

\paragraph{Warum kein Token im localStorage.} Eine naheliegende Alternative ist, das Token nach der Anmeldung im \texttt{localStorage} abzulegen und nach dem Verbindungsaufbau als erste WebSocket-Nachricht zu senden. Auf den \texttt{localStorage} kann jedoch jedes Skript zugreifen, das auf der Seite ausgeführt wird. Gelingt einem Angreifer \gls{xss}, etwa über eine Lücke in einer eingebundenen Bibliothek, kann er das Token auslesen und bis zum Ablauf verwenden. Deshalb wird davon abgeraten, Sitzungs-Tokens dort zu speichern \X[OWASP HTML5 Security Cheat Sheet, Local Storage]. Ein Token als Query-Parameter in der \gls{url} der Verbindung scheidet ebenfalls aus, weil \gls{url}s in Server-Logs und im Browserverlauf gespeichert werden \X[OWASP Session Management Cheat Sheet].

\paragraph{Session-Cookies.} Das System verwendet daher das für Browser-Sitzungen empfohlene Verfahren: Die Tokens werden ausschließlich in Cookies gespeichert, die der Server setzt \X[OWASP Session Management Cheat Sheet; RFC 6265]. Durch \texttt{HttpOnly} sind die Cookies für JavaScript nicht zugänglich, sodass auch ein eingeschleustes Skript die Tokens nicht auslesen kann. Der Browser sendet Cookies bei Anfragen an denselben Ursprung automatisch mit, auch beim WebSocket-Handshake. Der Client-Code berührt die Tokens damit an keiner Stelle.

\paragraph{Access- und Refresh-Token.} Ein einzelnes, lange gültiges Token hätte zwei Nachteile: Wird es gestohlen, ist es lange nutzbar, und der Server kann es nicht widerrufen, da er über ein \gls{jwt} nichts speichert. Ein sehr kurz gültiges Token allein würde dagegen bedeuten, dass sich Nutzer nach jedem Neuladen der Seite erneut anmelden müssen. Das übliche Vorgehen kombiniert deshalb zwei Tokens \X[OWASP JWT Cheat Sheet; RFC 9700 OAuth 2.0 Security BCP, Refresh Token Rotation].

\begin{itemize}
  \item Das \textbf{Access-Token} ist ein \gls{jwt} mit der Spieler-ID und der ID der Sitzung. Es ist 15 Minuten gültig und wird beim WebSocket-Handshake geprüft. Für die Prüfung genügt die Signatur, eine Datenbankabfrage ist nicht nötig.
  \item Das \textbf{Refresh-Token} ist eine zufällige Zeichenkette ohne eigene Bedeutung. Es ist 14 Tage gültig und dient nur dazu, ein neues Access-Token zu erhalten. Der Server speichert es als SHA-256-Hash in der Datenbank. Dadurch kann er es jederzeit sperren. Wer die Datenbank liest, erhält keine verwendbaren Tokens.
\end{itemize}

Das Attribut \texttt{Path=/auth} sorgt dafür, dass der Browser das Refresh-Token nur an \texttt{/auth/refresh} und \texttt{/auth/logout} sendet, nicht an alle anderen Pfade. \texttt{SameSite=Strict} verhindert, dass ein Cookie bei Anfragen mitgesendet wird, die von einer fremden Webseite ausgelöst werden. \texttt{Secure} beschränkt die Übertragung auf \gls{tls} und wird gesetzt, wenn \texttt{NODE\_ENV=production} gilt.

\paragraph{Erneuerung und Rotation.} Bei jedem Aufruf von \texttt{/auth/refresh} sucht der Server den Hash des vorgelegten Refresh-Tokens in der Datenbank. Ist der Eintrag gültig, markiert er ihn als verbraucht und legt ein neues Refresh-Token derselben Sitzung an. Beide Schritte laufen im Datenbank-Worker in einer Transaktion. Dort wird auch das neue Token zufällig erzeugt, in der Datenbank landet nur sein Hash. Erst wenn der Worker das neue Token zurückgegeben hat, setzt der Server beide Cookies neu. Das neue Refresh-Token ist wieder 14 Tage gültig. Wer die Anwendung regelmäßig nutzt, bleibt dadurch angemeldet, wer 14 Tage inaktiv ist, muss sich neu anmelden. Wird ein bereits verbrauchtes Refresh-Token erneut vorgelegt, besitzt vermutlich jemand eine Kopie. Der Server sperrt dann alle Tokens dieser Sitzung, sodass sich Angreifer und Nutzer neu anmelden müssen \X[RFC 9700, Refresh Token Rotation]. Eine Ausnahme bilden Tokens, die erst vor weniger als 30 Sekunden ersetzt wurden. Öffnet ein Nutzer mehrere Tabs gleichzeitig, können zwei Anfragen mit demselben Token kurz nacheinander eintreffen. Die zweite wird dann ohne Sperre mit 200 beantwortet, weil die Antwort auf die erste die neuen Cookies bereits gesetzt hat.

\paragraph{Ablauf bei abgelaufenem Access-Token.} Abbildung~\ref{fig:refresh} zeigt den Ablauf, wenn ein Nutzer die Seite nach mehr als 15 Minuten neu lädt. Der Handshake wird mit 401 abgelehnt. Der Client ruft daraufhin einmal \texttt{/auth/refresh} auf und baut die Verbindung mit dem neuen Access-Token erneut auf. Der Nutzer bemerkt davon nichts. Nur wenn auch die Erneuerung scheitert, erscheint die Anmeldeseite. Eingaben gehen dabei nicht verloren, weil dieser Ablauf vor dem Aufbau der Verbindung stattfindet. Eine offene Verbindung ist davon nicht betroffen: Das Access-Token wird nur beim Handshake geprüft, und die Verbindung bleibt nach dessen Ablauf bestehen.

\begin{figure}[ht]
\centering
\begin{tikzpicture}[font=\sffamily\scriptsize, >=Stealth,
  head/.style={draw, rounded corners=2pt, fill=black!4, minimum width=2.4cm, minimum height=0.7cm, font=\sffamily\small},
  every node/.append style={inner sep=1.5pt}]
  \node[head] at (0,0) {Browser};
  \node[head] at (6,0) {Server};
  \node[head] at (11,0) {Datenbank-Worker};
  \foreach \x in {0,6,11} \draw[black!40] (\x,-0.35) -- (\x,-7.4);
  \draw[->] (0,-0.9) -- node[above, fill=white] {Handshake \texttt{/ws}, Cookie: abgelaufenes Access-Token} (6,-0.9);
  \draw[->, dashed] (6,-1.6) -- node[above, fill=white] {\texttt{401 Unauthorized}} (0,-1.6);
  \draw[->] (0,-2.4) -- node[above, fill=white] {\texttt{POST /auth/refresh}, Cookie: Refresh-Token} (6,-2.4);
  \draw[->] (6,-3.0) -- node[above, fill=white] {Refresh-Token rotieren (eine Transaktion)} (11,-3.0);
  \draw[->, dashed] (11,-3.6) -- node[above, fill=white] {Ergebnis: Spieler, Sitzung, neues Refresh-Token} (6,-3.6);
  \draw[->, dashed] (6,-4.3) -- node[above, fill=white] {\texttt{200}, \texttt{Set-Cookie}: neues Access- und Refresh-Token} (0,-4.3);
  \draw[->] (0,-5.1) -- node[above, fill=white] {Handshake \texttt{/ws}, Cookie: neues Access-Token} (6,-5.1);
  \draw[->, dashed] (6,-5.8) -- node[above, fill=white] {\texttt{101 Switching Protocols}} (0,-5.8);
  \draw[->, dashed, thick] (6,-6.6) -- node[above, fill=white] {\texttt{welcome}} (0,-6.6);
\end{tikzpicture}
\caption{Erneuerung der Sitzung, wenn das Access-Token beim Verbindungsaufbau abgelaufen ist}
\label{fig:refresh}
\end{figure}

\paragraph{Abmeldung.} \texttt{/auth/logout} sperrt alle Refresh-Tokens der Sitzung in der Datenbank und löscht beide Cookies. Da eine offene Verbindung nach Ablauf des Access-Tokens bestehen bleibt, muss der Server sie aktiv beenden. Dafür speichert er beim Verbindungsaufbau die Sitzungs-ID aus dem Access-Token. Bei der Abmeldung schließt er die Verbindung dieser Sitzung mit dem Statuscode \texttt{4000} (\enquote{Logged out}). Damit ist eine Sitzung sofort beendet und kann nicht mehr erneuert werden.

\paragraph{Prüfung beim Handshake.} Der Server prüft das Access-Token, bevor überhaupt eine WebSocket-Verbindung entsteht. Er wertet dazu den \gls{http}-Upgrade-Request aus (Abschnitt~\ref{sec:ws-server}). Fehlt ein gültiges, nicht abgelaufenes Access-Token, beantwortet er den Handshake mit \texttt{401 Unauthorized} und schließt die \gls{tcp}-Verbindung. Ist das Token gültig, wird die Verbindung aufgebaut und fest mit dem Spieler verknüpft. Der Server sendet dann sofort die \texttt{welcome}-Nachricht. Eine Anmeldenachricht des Clients gibt es nicht. Bei allen weiteren Nachrichten ermittelt der Server die Identität ausschließlich aus dieser Verknüpfung, nie aus dem Inhalt der Nachricht. Ein Client kann sich daher nicht als ein anderer Spieler ausgeben, indem er eine fremde ID in eine Nachricht schreibt.

\paragraph{Prüfung des Ursprungs.} Weil der Browser das Cookie automatisch mitsendet, könnte eine fremde Webseite im Browser eines angemeldeten Nutzers eine WebSocket-Verbindung zum Server öffnen (Cross-Site WebSocket Hijacking). Die Same-Origin-Policy, die das bei gewöhnlichen \gls{http}-Anfragen verhindert, gilt für WebSocket-Handshakes nicht. Der Server prüft deshalb den Header \texttt{Origin}, den der Browser setzt und den Skripte nicht verändern können \X[RFC 6455, Abschnitt 10.2]. Er akzeptiert nur Ursprünge aus der Umgebungsvariable \texttt{ALLOWED\_ORIGINS} und lehnt alle anderen mit \texttt{403 Forbidden} ab. Ist die Variable nicht gesetzt, gelten die Adressen der lokalen Entwicklungsumgebung. Zusätzlich verhindert \texttt{SameSite=Strict}, dass das Cookie bei Anfragen von fremden Seiten mitgesendet wird.


Abbildung~\ref{fig:devtools-headers} zeigt den Handshake im laufenden System: Der Browser sendet das Access-Token-Cookie \texttt{session} und den Header \texttt{Origin} mit, und der Server antwortet mit \texttt{101 Switching Protocols}.

\paragraph{Mehrfach-Login.} Meldet sich ein Spieler an, der bereits verbunden ist, etwa in einem zweiten Browser, schließt der Server die ältere Verbindung mit dem Statuscode \texttt{1008} und dem Grund \enquote{Duplicate login}. Der Client erkennt diesen Code, setzt seinen Zustand zurück, wechselt zur Anmeldeseite und zeigt den Hinweis \enquote{Logged in on another device}. Damit besitzt jeder Spieler zu jedem Zeitpunkt höchstens eine aktive Verbindung.

\begin{figure}[ht]
  \centering
  \screenshot[\linewidth]{devtools_headers.png}
  \caption{Opening Handshake der Verbindung zu \texttt{/ws} in den Chrome-DevTools: Der Request enthält das Access-Token-Cookie und den \texttt{Origin}-Header, der Server antwortet mit 101}
  \label{fig:devtools-headers}
\end{figure}

\section{Datenhaltung}
\label{sec:daten}

Das System unterscheidet zwischen kurzlebigen und dauerhaften Daten. Kurzlebige Daten existieren nur, solange der Server läuft, und liegen in drei \texttt{Map}-Objekten im Speicher (Modul \texttt{state.js}):

\begin{itemize}
  \item \texttt{lobbies}: alle offenen Lobbys mit Spielern, Bereit-Status, Passwort und Chatverlauf,
  \item \texttt{activePlayers}: alle verbundenen Spieler mit ihrer WebSocket-Verbindung,
  \item \texttt{pendingInvites}: offene Lobby-Einladungen je Spieler.
\end{itemize}

Diese Daten beziehen sich auf aktuell verbundene Spieler und verlieren ihre Bedeutung, sobald die Verbindungen enden. Eine Lobby ohne Spieler wird gelöscht.

Dauerhaft in einer SQLite-Datenbank gespeichert werden Spielerkonten, Freundschaftsanfragen, Freundschaften, private Nachrichten und die Refresh-Tokens der Sitzungen. Diese Daten müssen einen Neustart des Servers überstehen. Das Schema ist in Abschnitt~\ref{sec:sqlite} beschrieben.

Da Node.js alle Handler nacheinander im Haupt-Thread ausführt (Abschnitt~\ref{sec:express}), können die \texttt{Map}-Objekte ohne Sperren gelesen und verändert werden. Die Datenbank liegt in einem eigenen Worker Thread (Abschnitt~\ref{sec:worker}), der keinen Speicher mit dem Haupt-Thread teilt.
